%% kapa: exam=1 round=32 year=2021 period=2 form=A subject=appraisal_law subject_ko=감정평가관계법규 seq=25 number=25
%% source: Q-Net 감정평가사 기출문제 (cst00309) artlSeq=5208397 file=extracted/1008/5208397/A형/2021년도 제32회 감정평가사 1차 2교시 A형.pdf
%% conversion: kordoc 4.24.51 -> tools/gen1.py (auto); review=proofread-sonnet; flags=none
\begin{question}{25}
건축법령상 이행강제금에 관한 설명으로 옳은 것은?
\begin{choices}
  \item 이행강제금은 건축신고 대상 건축물에 대하여 부과할 수 없다.
  \item 이행강제금의 징수절차는 「지방세법」을 준용한다.
  \item 허가권자는 이행강제금을 부과하기 전에 이행강제금을 부과·징수한다는 뜻을 미리 문서로써 계고하여야 한다.
  \item 허가권자는 위반 건축물에 대한 시정명령을 받은 자가 이를 이행하면 이미 부과된 이행강제금의 징수를 즉시 중지하여야 한다.
  \item 허가권자는 최초의 시정명령이 있었던 날을 기준으로 하여 1년에 5회 이내의 범위에서 그 시정명령이 이행될 때까지 반복하여 이행강제금을 부과·징수할 수 있다.
\end{choices}
\end{question}
\answer{3}
